% ---
% title: Calculate $log \pi_\theta$ in code
% date: 2025/11/17
% tag: Diffusion
% ---
\definecolor{mdc1}{HTML}{FFFFFF}
\definecolor{mdc2}{HTML}{FF0000}
在扩散模型（Diffusion  Model）的上下文中，\(log\Pi_{\theta}\) 指的是\colorbox{mdc2}{\textcolor{mdc1}{生成整张图片过程的对数概率}}。它不是一次性计算出来的，而是反向扩散（去噪）过程中\colorbox{mdc2}{\textcolor{mdc1}{每一步的对数概率的累加和}}。

这个计算主要发生在 \verb|flow_grpo/diffusers_patch/pipeline_with_logprob.py| 文件中。

\section{核心思想}

反向去噪过程可以看作是一个马尔可夫链：

\[
p(x_0) = \int p(x_0|x_1) p(x_1|x_2) \dots p(x_{T-1}|x_T) p(x_T) dx_{1 \dots T}
\]

其中，\(p(x_T)\) 通常是一个标准正态分布，而每一步 \(p(x_{t-1}|x_t, \text{prompt})\) 都被建模成一个高斯分布。

因此，生成一张完整图片 \(x_0\) 的总对数概率 \(log P(x_0)\) 就是从 \(t=T\)到 \(t=1\) 每一步的对数概率 \(log p(x_{t-1}|x_t)\) 的总和。

\[
\log \Pi = \log p(x_0) \approx \sum_{t=1}^{T} \log p(x_{t-1} | x_t, \text{prompt})
\]

\section{具体实现}

在您项目中的 \verb|pipeline_with_logprob| 函数内，这个计算过程大致如下：

\begin{enumerate}
  \item 初始化: 在进入主要的去噪循环之前，会初始化一个 \verb|log_prob| 累加器，通常是一个全零的张量，形状与批次大小 (\verb|batch size|)
  相同。
  \item 去噪循环 (Denoising Loop): 代码会遍历所有的时间步 \(t\) (从 \(T\) 到 \(1\))。在循环的每一步中：
  
  \begin{itemize}
    \item 模型预测: 扩散模型（在您的代码中是 transformer 或 unet）以当前的噪声图像 latents (\(x_t\)) 和时间步 \verb|t|作为输入，预测出在当前步骤应该被移除的噪声 \verb|noise_pred|。
    \item 计算上一步的分布参数: 使用 \verb|noise_pred| 和当前的 latents (\(x_t\))，以及调度器 (\verb|scheduler|)的数学公式，可以计算出\textbf{上一步} latents\_prev (\(x_{t-1}\))的高斯分布的均值（mean）和方差（variance）/对数方差（\verb|log_variance|）。
    \item 获取实际的 \verb|latents_prev|: 调度器的 step 函数会实际生成上一步的噪声图像 \verb|latents_prev|。
    \item 计算单步 \verb|log_prob|: 现在我们有了 \(p(x_{t-1}|x_t)\) 的高斯分布参数（均值和方差），也知道了实际生成的样本点\verb|latents_prev|。于是，就可以使用高斯分布的概率密度函数 (PDF) 来计算 \verb|latents_prev| 在这个分布下的对数概率。
    
    高斯分布的对数概率公式为：
    
    \[
    \log \mathcal{N}(x; \mu, \sigma^2) = -\frac{1}{2} \log(2\pi\sigma^2) - \frac{(x - \mu)^2}{2\sigma^2}
    \]
    
    代码里通常会计算每个像素的 \verb|log_prob|，然后求和或求平均。
    \item 累加: 将这一步计算出的 \verb|log_prob| 加到总的累加器上。
  \end{itemize}
  \item 返回结果: 循环结束后，\verb|log_prob| 累加器里就存储了生成这批次中每张图片的总对数概率。
\end{enumerate}

\section{总结}

简单来说，\(log \Pi\) 的计算过程是：

\begin{enumerate}
  \item 在每个去噪步骤 \(t\)，模型会预测一个用于描述 \(x_{t-1}\) 分布的“目标”（通常是均值）。
  \item 将实际生成的 \(x_{t-1}\) 与这个分布进行比较，计算其在该分布下的对数概率。
  \item 将所有步骤的对数概率相加，得到最终生成整张图片的\( log\Pi\)。
\end{enumerate}

这个值在强化学习微调中至关重要，因为它代表了策略（扩散模型）在给定状态（prompt）下，采取某个动作（生成特定图片）的概率 ，是策略梯度算法 (如 PPO) 所必需的核心要素。
